%% paper_additions.tex
%% New sections to integrate into paper_shared_body.tex
%%
%% Integration points:
%%   Section A → append to \section{Dataset and Experimental Setup}
%%   Section B → new \section, insert before \section{Results}
%%   Section C → append to \section{Results}
%%   Section D → new \section after results, before conclusion
%%   Section E → replace bullet-list contributions in \section{Introduction}


%% ═══════════════════════════════════════════════════════════════════════════
%% SECTION A  –  Graph Construction Variants
%% (Append to sec:dataset, after subsection{Evaluation Protocol})
%% ═══════════════════════════════════════════════════════════════════════════

\subsection{Graph Construction Variants}
\label{sec:graph_construction}

The Elliptic transaction graph is sparse by design: each transaction has
on average 2.3 neighbours (Table~\ref{tab:graph_stats}), and 49 connected
components correspond precisely to the 49 temporal slices, with no edges
crossing slice boundaries. Whether these transaction-flow edges are the most
informative graph structure for fraud detection is not obvious. A node's
direct BTC counterparties may be less diagnostic than nodes that share
similar behaviour, or nodes active in adjacent time windows.

We therefore construct four graph variants from the same node features and
evaluate each under identical training conditions.

\textbf{Similarity graph.} For each time step independently, we compute the
cosine similarity between every pair of node feature vectors and add an
undirected edge wherever the similarity exceeds 0.92. Operating within
timesteps keeps the computation tractable (roughly 4,000 nodes per slice)
while ensuring comparisons are made between contemporaneous transactions.
This yields 3.6M edges and a mean degree of 17.7.

\textbf{Feature k-NN graph.} Each node is connected to its five nearest
neighbours in feature space (L2 distance) within the same timestep. The
resulting graph is denser than the original but sparser than the similarity
graph: 1.0M edges, mean degree 4.9.

\textbf{Temporal graph.} For adjacent timesteps $t$ and $t{+}1$, each node
in $t$ is linked to its three nearest feature-space neighbours in $t{+}1$.
This encodes temporal continuity: transactions that resemble each other
across consecutive slices are connected even when no direct BTC flow exists
between them. The graph has 1.2M edges, mean degree 5.9, and no within-step
edges.

\textbf{Augmented graph.} The original transaction edges combined with all
similarity edges, yielding 4.1M edges and mean degree 19.9.

\begin{table}[!htbp]
\centering
\caption{Topology of each graph construction variant. Connected-components
  values are omitted for variants with more than 800K edges (marked ---)
  because the NetworkX computation becomes prohibitively slow on the
  full 203K-node graph. Clustering is estimated over a 3,000-node sample.}
\label{tab:graph_stats}
\small
\begin{tabular}{lrrrrr}
\toprule
\textbf{Graph type} & \textbf{Edges} & \textbf{Mean deg.}
  & \textbf{Max deg.} & \textbf{Clustering} & \textbf{CC} \\
\midrule
Original    &    468,710 &  2.30 & 473 & 0.000 & 49 \\
Similarity  &  3,611,346 & 17.72 & 482 & 0.045 & ---\\
Feature k-NN & 1,007,816 &  4.95 &  25 & 0.000 & ---\\
Temporal    &  1,207,890 &  5.93 & 331 & 0.000 & ---\\
Augmented   &  4,057,870 & 19.91 & 484 & 0.045 & ---\\
\bottomrule
\end{tabular}
\end{table}

Several structural properties are worth noting before reporting performance.
The original graph has clustering coefficient zero across the sampled nodes:
no two neighbours of any node are themselves connected. Feature-based graphs
behave differently: in the similarity and augmented cases, clustering rises to
0.045 because shared feature similarity is transitive. The high variance in
max-degree (473 in the original, 25 in feature k-NN) also matters: the
original graph contains high-degree hub transactions that are common in
Bitcoin mixing services, a structural signature that feature-based graphs
cannot replicate.


%% ═══════════════════════════════════════════════════════════════════════════
%% SECTION B  –  Comparison Against Classical Baselines
%% (Insert as new \section before \section{Results})
%% ═══════════════════════════════════════════════════════════════════════════

\section{Comparison Against Classical Baselines}
\label{sec:classical_comparison}

Before examining graph-construction effects on GNN performance, it is worth
establishing what classical ML achieves on the same features without any
graph structure.

\subsection{Baseline models}

We train three models on the 165 node features directly, with no message
passing or neighbourhood aggregation.

\textbf{Logistic Regression} with class-balanced weights and L-BFGS solver
serves as the linear lower bound. It is computationally negligible and
interpretable; its coefficients are the simplest possible summary of which
features matter.

\textbf{Random Forest} with 300 trees, balanced class weights, and minimum
leaf size 2. This is the strongest classical result in the Elliptic
literature; Weber et al.\ \citep{weber2019} reported F1\,=\,0.64 under
transductive evaluation. We report it under the same strict temporal split
used for GNNs.

\textbf{XGBoost} with 300 boosting rounds, max depth 6, learning rate 0.05,
and scale\_pos\_weight equal to the negative-to-positive ratio (7.6). Tree
construction uses the histogram method.

Results are in Table~\ref{tab:classical_vs_gnn}.

\begin{table}[!htbp]
\centering
\caption{Classical ML vs. GNN-SAGE on the original graph (single run).
  The GNN numbers here are from a single training run; the best GNN result
  in the multi-seed benchmark (GraphSAGE + weighted loss, 5 seeds) reaches
  F1\,=\,0.781\,$\pm$\,0.007, closing most but not all of the gap with RF.
  AUC scores are closer than F1 scores, suggesting the GNN probability
  ranking is reasonable but its decision threshold is misaligned.}
\label{tab:classical_vs_gnn}
\small
\begin{tabular}{lcccc}
\toprule
\textbf{Model} & \textbf{F1} & \textbf{Precision} & \textbf{Recall}
  & \textbf{AUC} \\
\midrule
Logistic Regression & 0.314 & 0.192 & 0.869 & 0.887 \\
GNN-SAGE (original graph) & 0.290 & 0.179 & 0.764 & 0.845 \\
XGBoost          & 0.775 & 0.819 & 0.735 & 0.931 \\
GNN-SAGE + RF (hybrid) & 0.558 & 0.494 & 0.640 & 0.830 \\
Random Forest    & \textbf{0.820} & 0.963 & 0.715 & \textbf{0.933} \\
\bottomrule
\end{tabular}
\end{table}

\subsection{What the numbers actually say}

Random Forest achieves F1\,=\,0.820 and AUC\,=\,0.933 with no graph
structure at all. The single-run GNN on the original transaction graph
achieves F1\,=\,0.290, roughly 2.8$\times$ lower. Even with proper
multi-seed tuning (F1\,=\,0.781), the best GNN still falls short of RF by
0.039 F1 points. Two things explain this.

First, the Elliptic feature set already encodes a form of manual message
passing. Features 94--164 are precomputed neighbourhood aggregates supplied
by the dataset creators; they give any feature-only model access to first-hop
graph information for free. The GNN is therefore not competing against a
completely graph-blind baseline.

Second, the decision-threshold problem is visible in the AUC numbers. The
temporal GNN achieves AUC\,=\,0.895, only 0.038 below RF's AUC\,=\,0.933.
The GNN probability ordering is reasonable; the issue is that its F1-optimal
threshold shifts over time as the fraud rate changes, and a fixed default
threshold of 0.5 performs poorly during low-fraud-rate test windows.

\subsection{The hybrid model does not help}

Passing GNN-SAGE embeddings (256-dimensional, extracted before the final
linear layer) into a Random Forest gives F1\,=\,0.558 — considerably worse
than RF on raw features. This is not the expected outcome for a hybrid model.
Two likely explanations: the GNN embeddings are trained under class imbalance
with weighted loss, which warps the embedding space away from the geometry RF
exploits; and the 256-dimensional embeddings from a single-run GNN carry less
per-feature information than the 165 handcrafted Elliptic features that RF
was designed to split on. In short, GNN embeddings add noise, not signal,
when the base feature set is already strong.


%% ═══════════════════════════════════════════════════════════════════════════
%% SECTION C  –  Graph Construction Results
%% (Append to \section{Results}, after the per-timestep subsection)
%% ═══════════════════════════════════════════════════════════════════════════

\subsection{Effect of graph construction on GNN performance}
\label{sec:graph_construction_results}

Table~\ref{tab:graph_results} shows GraphSAGE performance across the five
graph variants under identical training conditions (single run, 300 epochs,
weighted loss).

\begin{table}[!htbp]
\centering
\caption{GraphSAGE performance by graph construction (single run per
  variant). The original transaction graph gives the worst GNN F1 of
  all five variants. A temporal cross-timestep graph nearly doubles F1
  at the cost of discarding the actual Bitcoin transaction structure.}
\label{tab:graph_results}
\small
\begin{tabular}{lcccc}
\toprule
\textbf{Graph} & \textbf{F1} & \textbf{Precision} & \textbf{Recall}
  & \textbf{AUC} \\
\midrule
Original (transaction)    & 0.290 & 0.179 & 0.764 & 0.845 \\
Similarity (cosine)       & 0.362 & 0.246 & 0.680 & 0.842 \\
Augmented (orig + simil.) & 0.385 & 0.259 & 0.751 & 0.867 \\
Feature k-NN              & 0.453 & 0.336 & 0.696 & 0.862 \\
Temporal (cross-step)     & \textbf{0.549} & 0.456 & 0.690 & \textbf{0.895} \\
\bottomrule
\end{tabular}
\end{table}

The original transaction graph performs worst despite being the semantically
correct representation of the problem. The temporal graph performs best
despite containing no transaction-flow information at all — its edges are
defined entirely by feature similarity between adjacent time steps. This
inversion is explained by the average degree numbers in
Table~\ref{tab:graph_stats}: the original graph's 2.3 mean degree gives each
node only two neighbours for message passing, providing a thin information
channel. The temporal graph's 5.9 mean degree more than doubles the available
neighbourhood context.

The ablation results confirm this interpretation. Removing all edges from the
original graph actually improves F1 from 0.290 to 0.341 (Table~\ref{tab:ablation}).
The original transaction edges are marginally harmful rather than helpful: their
sparsity forces the GNN to aggregate over a narrow neighbourhood whose content
is already partially captured by the precomputed aggregate features (columns
94--164). The constructed graphs provide richer context, but only the temporal
graph provides enough context to matter.

\begin{table}[!htbp]
\centering
\caption{Ablation results. ``Structural features only'' uses five
  hand-computed structural attributes (degree, log-degree, normalised
  timestep, in-degree, out-degree) as the sole node features. The
  complete feature set contains 165 dimensions. The results show that
  discriminative power is almost entirely carried by node content
  features, not graph structure.}
\label{tab:ablation}
\small
\begin{tabular}{llcc}
\toprule
\textbf{Experiment} & \textbf{What was removed} & \textbf{F1} & \textbf{AUC} \\
\midrule
Full model (GNN + original graph) & Nothing       & 0.290 & 0.845 \\
No edges (MLP over features)      & Graph structure & 0.341 & 0.866 \\
Structural features only          & Node content  & 0.000 & 0.598 \\
No edges + structural features    & Graph + content & 0.000 & 0.589 \\
\bottomrule
\end{tabular}
\end{table}

The structural-features-only result is unambiguous: degree, timestep, and
in/out-degree carry zero discriminative power for fraud detection on this
dataset. The model predicts the majority class for every node, achieving
accuracy of 35.6\,\% but F1 of exactly zero. Fraud detection on Elliptic
is a content problem, not a topology problem. Graph structure can assist
when the topology is informative (as in the temporal variant), but it cannot
compensate for the absence of content features.


%% ═══════════════════════════════════════════════════════════════════════════
%% SECTION D  –  Key Insights
%% (New \section after Results, before Conclusion)
%% ═══════════════════════════════════════════════════════════════════════════

\section{Key Insights}
\label{sec:insights}

The following observations are drawn from the full set of experiments: the
multi-seed GNN benchmark, the graph construction study, the classical
baseline comparison, the hybrid model evaluation, and the ablation.
Each addresses a question that the prior Elliptic literature leaves open.

\subsection{When graph structure helps}

Graph structure provides measurable benefit in two conditions.

The first is when the constructed neighbourhood is rich enough to provide
context that is not already redundant with precomputed features. The temporal
graph, with mean degree 5.9 and cross-timestep connections, achieves the
highest GNN AUC (0.895) and F1 (0.549) of all five graph variants. The
original transaction graph, with mean degree 2.3 and clustering coefficient
zero, provides the least benefit. More neighbours means more context,
and constructed graphs tuned for feature similarity deliver that context
even when the underlying transaction graph does not.

The second is during stable fraud-rate windows. The per-timestep analysis
(Section~\ref{sec:results}) shows that models — GNN and classical alike —
are meaningfully better during test steps 35--42, when the fraud rate
stays within $2\times$ the training distribution. During this window,
GAT with graph augmentation achieves F1\,=\,0.789, a genuine advantage
over a feature-only baseline.

\subsection{When graph structure does not help}

The original transaction graph edges, taken at face value, actively hurt
performance relative to no edges at all. The ablation removes all edges
and gains 0.051 F1 points. The graph augmentation strategy used in the
multi-seed benchmark achieves its advantage by adding synthetic ego-graph
subgraphs around minority nodes — it is an oversampling strategy with
graph structure bolted on, not pure structural reasoning.

The hybrid model result is telling: Random Forest on raw features beats
Random Forest on GNN embeddings by 0.262 F1 points. The GNN embeddings
do not compress the feature space into something more useful; they compress
it into something less useful. A possible explanation is that the 165
handcrafted features were engineered by domain experts who understood what
discriminates Bitcoin fraud from legitimate activity; the GNN re-encodes
those features through a generic message-passing lens that has no access to
that prior knowledge.

The temporal collapse (steps 43--49) is immune to graph structure. Both
feature-only and graph-based models fail when the fraud rate drops to 0.3\,\%.
This is a prior-shift problem, not an architecture problem. No graph
construction variant recovers from it.

\subsection{What the AUC gap tells us}

The best GNN (GAT + graph augmentation, AUC\,=\,0.919) and the best classical
model (Random Forest, AUC\,=\,0.933) are separated by 0.014 AUC units.
The F1 gap between the same models, under a fixed decision threshold, is
much larger. This pattern is consistent with a calibration mismatch that
worsens as fraud rates drift from training conditions. The GNN's probability
orderings are reasonable; its decision boundary is not.
This is a calibration problem, not an expressivity problem, and threshold
retuning on a held-out validation window should close most of the F1 gap
during stable periods.

\subsection{On feature dominance}

Features 40--60 in the local feature block appear in the top-ranked
positions for both RF Gini importance and GNN gradient saliency. The
feature indices with highest GNN saliency (f52, f51, f80) overlap with
the top RF features (f54, f46, f52), with Spearman rank correlation of
approximately 0.71 across the top-15 features. Both models agree on which
transaction properties matter most, even though they discover this agreement
through completely different mechanisms. This convergence is modest evidence
that the important features reflect genuine fraud structure rather than
statistical artifacts specific to any one model family.

\subsection{Practical takeaway}

For a production fraud detector on a dataset with these properties, the
most defensible choice is still a well-tuned ensemble (Random Forest or
XGBoost) applied directly to the full feature set. The multi-seed
GraphSAGE result (F1\,=\,0.781) closes the gap with RF substantially,
but the 0.039 F1-point deficit comes with considerably more engineering
overhead and a failure mode during fraud-rate shifts that classical models
share. The main practical case for GNNs here is the temporal graph
construction result: if the transaction graph is replaced by a richer
constructed graph capturing behavioural similarity over time, GNN
performance improves by 89\,\% relative (F1 0.290 $\to$ 0.549). Whether
that is worth the construction overhead depends on the deployment context.
On a live Bitcoin network where transaction edges are a direct artifact of
the ledger and feature extraction is the expensive step, the temporal graph
construction requires computing and indexing per-timestep nearest neighbours
— feasible in batch, nontrivial in real time.


%% ═══════════════════════════════════════════════════════════════════════════
%% SECTION E  –  Updated contribution list
%% (Replace the bullet-list contributions in \section{Introduction})
%% ═══════════════════════════════════════════════════════════════════════════

%% Replace the existing \begin{itemize}...\end{itemize} contributions block
%% in the introduction with the following:

\begin{itemize}[leftmargin=*]
  \item A statistically validated benchmark across 4 architectures
    $\times$ 3 imbalance strategies $\times$ 5 random seeds under
    inductive mini-batch evaluation, establishing the single-number
    performance ceiling for GNNs on this dataset.

  \item A systematic comparison of five graph construction strategies
    (original transaction graph, cosine similarity, feature k-NN,
    cross-timestep temporal, and augmented). The temporal graph
    construction improves GNN F1 by 89\,\% relative over the original
    transaction graph — the most substantial intervention tested.

  \item A direct comparison against Random Forest and XGBoost under the
    same temporal split. Classical methods outperform single-run GNNs;
    the gap narrows under multi-seed averaging but does not close.

  \item An ablation decomposing GNN performance into contributions from
    graph structure (edges) and node content (features), showing that
    the original transaction graph is marginally harmful and that node
    content features carry virtually all discriminative power.

  \item A hybrid model evaluation (GNN embeddings $\to$ Random Forest)
    showing that GNN representations do not improve over raw features
    when the feature set is already expressive, with a 0.262 F1-point
    decline relative to RF on raw features.

  \item Feature importance analysis across RF (Gini), XGBoost (gain),
    and GNN (input-gradient saliency), showing that models from
    different families converge on the same feature subspace.

  \item The first systematic per-timestep analysis on Elliptic,
    demonstrating that aggregate test-set F1 conceals near-complete
    failure in steps 43--49 and an impossibility result for post-hoc
    recalibration under those conditions.

  \item Code, checkpoints, and experiment scripts:
    \url{https://github.com/Saket-Maganti/gnn-fraud}.
\end{itemize}
